\documentclass[12pt,a4paper,oneside]{book}
\usepackage[utf8]{inputenc}
\usepackage[T1]{fontenc}
\usepackage[english,french]{babel}
\usepackage{geometry}
\usepackage[table]{xcolor}

\geometry{top=2.5cm,bottom=2.5cm,left=2.5cm,right=2.5cm}
\definecolor{bleuIA}{RGB}{37, 99, 235}

\begin{document}

\begin{titlepage}
\begin{center}
\vspace*{3cm}
{\Huge\bfseries\color{bleuIA}E-GRAPHISME}\\[0.3cm]
{\LARGE by ELECTRON}\\[2cm]
{\huge ENTERPRISE AI OPERATING SYSTEM}\\[0.3cm]
{\huge Architecture Technique Globale}\\[1cm]
{\Large Volume 16}\\[3cm]
{\Large Version 2026}
\end{center}
\end{titlepage}

\chapter{Vision}
Le système fonctionne comme un véritable OS d'entreprise.

\chapter{Architecture}
\begin{itemize}
\item Frontend React + Next.js
\item API Backend
\item Services IA
\item Base de données
\item Bus d'événements
\end{itemize}

\chapter{Microservices}
CRM, ERP, Boutique, Marketplace, Paiements, Formation.

\chapter{API Gateway}
Routage, authentification, limitation de débit.

\chapter{Identity Provider}
Gestion utilisateurs, rôles, permissions.

\chapter{AI Memory}
Mémoire pour les agents.

\chapter{Ollama Center}
Gestion des modèles IA.

\chapter{Prompt Management}
Bibliothèque centralisée.

\chapter{Security Zero Trust}
Vérification continue des identités.

\chapter{Scalabilité}
Ajout de nouveaux services.

\chapter{AI Model Router}
Sélection automatique du modèle IA adapté.

\chapter{RAG Center}
Gestion des connaissances.

\chapter{Infrastructure Manager}
Supervision technique.

\chapter{Plugin Center}
Système d'extensions.

\end{document}
